% SOURCE_STATEMENT: "Do the elliptic-cohomology partition functions for type II string theory coincide with partition functions of twelve-dimensional F-theory compactified on an elliptic curve? Test this proposed interpretation by computing suitable loop versions of the $E_8$ and Rarita–Schwinger indices."
% Preparation and final self-audit: CONTRIBUTING.md prompts C and D; source/audit artifacts retained in the local physics_statements_2026_10_11 batch.
% !TeX program = lualatex
\documentclass[11pt]{article}
\usepackage[margin=25mm]{geometry}
\usepackage{fontspec}
\setmainfont{Latin Modern Roman}
\usepackage{amsmath,amssymb,mathtools}
\usepackage{unicode-math}
\setmathfont{Latin Modern Math}
\usepackage{xurl,hyperref}
\hypersetup{colorlinks=true,urlcolor=blue}
\setcounter{secnumdepth}{0}
\setlength{\parindent}{0pt}
\setlength{\parskip}{5pt}
\setlength{\emergencystretch}{3em}
\newcommand{\sref}[2]{\hyperlink{source:#1:#2}{[S#2]}}
\newcommand{\cataloguescope}[1]{\paragraph{Recorded target.}#1}
\newcommand{\R}{\mathbb R}
\newcommand{\C}{\mathbb C}
\newcommand{\N}{\mathbb N}
\newcommand{\Z}{\mathbb Z}
\newcommand{\Q}{\mathbb Q}
\newcommand{\F}{\mathbb F}
\newcommand{\eps}{\varepsilon}
\begin{document}
% UnsolvedMath ID 30000244; source catalogue code OWR-832-003
\setcounter{section}{30000243}
\section{Physical Interpretation of Elliptic Partition Functions}\label{30000244}
{\small UnsolvedMath ID 30000244 · Mathematical Physics}
\subsection{Definitions and mathematical statement}

\paragraph{Shared notation.}$\mathbb N=\{1,2,\ldots\}$; the standard number fields have their usual meanings. All additional conventions are stated locally.


\paragraph{Definitions and conventions.} A partition function is the amplitude or generating function assigned to a physical theory on its spacetime background. Elliptic cohomology is a generalized cohomology theory associated with elliptic curves; the report uses it to refine the type IIA and IIB flux partition functions. The proposed F-theory interpretation introduces two extra dimensions represented by an elliptic curve. An index is the signed dimension of the kernel and cokernel of a suitable Dirac-type operator; the $E_8$ and Rarita--Schwinger indices refer to gauge and spin-$3/2$ contributions in the cited construction.

\paragraph{Source scope and qualifications.} The twelve-dimensional physical theory and its loop indices are proposed constructions. The report does not supply a fully axiomatized F-theory partition function with all backgrounds fixed; the target is an interpretation and its proposed tests, not an already defined equality for every conceivable background.

\paragraph{Catalogue mathematical target.}
\paragraph{Catalogue wording retained for traceability.}
Do the elliptic-cohomology partition functions for type II string theory coincide with partition functions of twelve-dimensional F-theory compactified on an elliptic curve? Test this proposed interpretation by computing suitable loop versions of the $E_8$ and Rarita–Schwinger indices.

\paragraph{Saved literature context.} The supplied catalogue's assessment dated 2026-08-21 records: The report states that the physical conjecture had not been verified; later index and partition-function work gives evidence and special calculations, not a general verification. The assessment is retained as dated provenance; it is not a new comprehensive literature or open-status certification.
\subsection{Short English statement}
Test whether elliptic-cohomology string partition functions arise by compactifying a twelve-dimensional F-theory, using the proposed loop-index calculations.
\subsection{Sources}
\begin{itemize}
\item[\hypertarget{source:30000244:1}{S1}] ULAM.AI and the UnsolvedMath Contributors, UnsolvedMath, record 30000244 (OWR-832-003). Catalogue curation under CC BY 4.0; definitions and English exposition added locally.
\url{https://huggingface.co/datasets/ulamai/UnsolvedMath}.
\item[\hypertarget{source:30000244:2}{S2}] Geometric Topology and Connections with Quantum Field Theory (2005), Oberwolfach Reports, DOI 10.4171/owr/2005/27, pp. 1505--1546; official report source and mathematical context.
\url{https://doi.org/10.4171/owr/2005/27}.
\item[\hypertarget{source:30000244:3}{S3}] Official EMS Press PDF of the same Oberwolfach report; PDF page 11.
\url{https://ems.press/content/serial-article-files/46001}.
\item[\hypertarget{source:30000244:4}{S4}] Kriz--Sati Type II modularity; reference in the saved dated catalogue literature review.
\url{https://arxiv.org/abs/hep-th/0501060}.
\end{itemize}
\label{end:30000244}
\end{document}
